% SPDX-License-Identifier: AGPL-3.0-or-later
% Commercial license available
% © Concepts 1996-2026 Miroslav Sotek. All rights reserved.
% © Code 2020-2026 Miroslav Sotek. All rights reserved.
% ORCID: 0009-0009-3560-0851
% Contact: www.anulum.li | protoscience@anulum.li
% scpn-quantum-control -- SCPN/FIM Hamiltonian manuscript draft

\documentclass[preprint,aps,physrev,superscriptaddress,floatfix,longbibliography]{revtex4-2}

\usepackage{amsmath,amssymb}
\usepackage{booktabs}
\usepackage{hyperref}
\usepackage{xcolor}
\usepackage{xurl}

\hypersetup{
    colorlinks=true,
    linkcolor=blue!60!black,
    citecolor=blue!60!black,
    urlcolor=blue!60!black,
}

\newcommand{\paperdoi}{Zenodo DOI: \href{https://doi.org/10.5281/zenodo.20382125}{10.5281/zenodo.20382125}}

\begin{document}
\emergencystretch=3em

\title{Fisher-information-inspired collective feedback in Kuramoto--XY
Hamiltonians: exact small-system structure and a hardware falsification}

\author{Miroslav \v{S}otek}
\email{protoscience@anulum.li}
\affiliation{ANULUM, Marbach SG, Switzerland}
\affiliation{ORCID: 0009-0009-3560-0851}

\date{\today\\\paperdoi}

\begin{abstract}
We introduce and characterise a collective Fisher-information-inspired
magnetisation-feedback term for Kuramoto--XY quantum simulations. The model
augments the heterogeneous XY Hamiltonian with a diagonal term proportional to
minus lambda times total magnetisation squared, normalised by the qubit count.
The term is diagonal in the
computational basis and therefore creates explicit magnetisation-sector energy
shifts without breaking the magnetisation conservation of the ideal XY
model. Artefact-generated exact diagonalisation on $n=4,6,8$ qubits shows that
increasing $\lambda$ lowers the ground energy, widens the spectrum, and
increases the spectral gap in the tested grid. For example, the exact gap
changes from $1.1317$ to $12.6060$ at $n=4$, from $0.4468$ to $13.7801$ at
$n=6$, and from $0.2827$ to $13.7173$ at $n=8$ when $\lambda$ is increased
from $0$ to $4$. The full-spectrum adjacent-gap mean shifts downward at
$n=6,8$ for larger $\lambda$, and the $n=4$ mean low-energy bipartition
entropy decreases from about $0.676$ bits at $\lambda\leq 1$ to about
$0.490$ bits at $\lambda=8$. A small $n=4$ VQE benchmark further shows that a
$K_{ij}$-informed ansatz gives lower median relative energy error than generic
\texttt{TwoLocal} and \texttt{EfficientSU2} baselines for
$\lambda\in\{0,1,4\}$. Finally, we executed the corresponding $n=4$ IBM Heron
r2 hardware tests on \texttt{ibm\_kingston}. A single-sample pilot was followed
by a repeated randomized run with 166 circuits and 339968 shots. The repeated
run is a negative validation result: at $\lambda=4$ relative to $\lambda=0$,
magnetisation leakage increased by $0.0896$ on average across matched
state/depth comparisons (Fisher combined $p=7.49\times10^{-55}$), while
exact-state retention decreased by $0.0797$ on average. A full 16-state
readout-confusion-matrix inversion on the same repeated run preserves this
negative conclusion: magnetisation leakage still increases by $0.0917$
($p=8.14\times10^{-55}$) and state retention still decreases by $0.0815$
($p=1.19\times10^{-48}$). We then executed a dedicated replication/ZNE lane
on \texttt{ibm\_marrakesh} and \texttt{ibm\_fez}. After full readout
mitigation and linear ZNE, the mean $\lambda=4$ minus $\lambda=0$
magnetisation-leakage deltas remained positive in the $n=4$ replication:
$+0.0599$ on Marrakesh and $+0.1937$ on Fez. Larger $n=6,8$ Fez/Marrakesh
replication/ZNE lanes were then executed. They preserve the falsification
boundary rather than promoting FIM protection: all four larger lanes show
positive mean raw scale-1 magnetisation-leakage deltas, with linear-ZNE deltas
of $0.4021$, $0.1480$, $0.2541$, and $0.4203$ for Fez $n=6$, Fez $n=8$,
Marrakesh $n=6$, and Marrakesh $n=8$, respectively. Thus the tested digital
Trotter implementation falsifies the simple hypothesis that this FIM term acts
as a hardware coherence-protection mechanism. The paper is deliberately
bounded: the artefacts support exact
small-system Hamiltonian findings and a backend/circuit-specific negative
hardware result on IBM Heron devices. They do not demonstrate
hardware coherence improvement, quantum advantage, strict many-body
localisation, or universal protection.
\end{abstract}

\maketitle

\section{Introduction}

Kuramoto synchronisation, XY spin Hamiltonians, and NISQ quantum simulation
provide a compact setting for studying how oscillator-network structure appears
in quantum dynamics~\cite{Kuramoto1975,Acebron2005,Barahona1982}. In the
standard mapping, heterogeneous oscillator frequencies and couplings are encoded
as local $Z$ fields and $XX+YY$ exchange terms. This construction preserves
total magnetisation because the XY interaction swaps excitations without
creating or destroying them.

This paper continues the SCPN Kuramoto--XY hardware programme, which first
established a bounded parity-sector and excitation-number correlated leakage
observation on IBM Heron hardware~\cite{SotekDLA2026} and then used
reduced-Pauli checks to localise correlator-level structure in the same
families~\cite{SotekPhase3ReducedPauli2026}. Here we test a distinct
Hamiltonian modification: a Fisher-information-inspired magnetisation-feedback
term. The contribution is a negative design rule: the direct digital Trotter
implementation of this FIM term is not a coherence-protection mechanism on the
tested superconducting hardware.

This manuscript studies a minimal extension of that model: a collective
magnetisation-feedback term inspired by Fisher-information and
self-consistency ideas. The term is
%
\begin{equation}
  H_{\mathrm{FIM}}(\lambda) = -\frac{\lambda}{n} M^2,
  \qquad
  M = \sum_{i=1}^{n} Z_i.
  \label{eq:fim}
\end{equation}
%
It is not introduced here as a validated error-correction mechanism. Rather,
the question is narrower and falsifiable: what does this term do to the exact
small-system energy landscape, sector structure, spectral diagnostics, and
low-energy entanglement, and does the most direct digital hardware
implementation show any corresponding leakage suppression?

The contribution is threefold. First, we define the FIM-augmented
Kuramoto--XY Hamiltonian and record the exact magnetisation-sector consequences
of Eq.~\eqref{eq:fim}. Second, we generate committed JSON/CSV artefacts for
spectra, level-spacing ratios, bipartition entropy, ideal sector-conservation
checks, and VQE ground-state scoring. Third, we report IBM Heron r2 raw-count
tests that falsify the simple claim that the tested FIM Trotter construction
improves hardware coherence on \texttt{ibm\_kingston}. This negative result is
part of the contribution: it separates a valid Hamiltonian-structure effect
from an unsupported NISQ robustness claim.

\section{Model}

The base Hamiltonian is the heterogeneous Kuramoto--XY model
%
\begin{equation}
  H_{XY} =
  -\sum_i \omega_i Z_i
  -\sum_{i<j} K_{ij}(X_iX_j+Y_iY_j),
  \label{eq:hxy}
\end{equation}
%
where $K_{ij}=K_{ji}$ is the oscillator coupling matrix and $\omega_i$ are
natural frequencies. The sign convention follows the implementation in
\texttt{scpn-quantum-control}. The augmented Hamiltonian is
%
\begin{equation}
  H(\lambda)=H_{XY}+H_{\mathrm{FIM}}(\lambda).
  \label{eq:total}
\end{equation}

Since $M$ is diagonal in the computational basis, $M^2$ assigns the same
energy shift to every state in a fixed magnetisation sector. For a basis state
with $k$ spin-down qubits, the magnetisation is
$M=n-2k$, and the FIM shift is
%
\begin{equation}
  \Delta E_M(\lambda)=-\lambda M^2/n.
  \label{eq:sector_shift}
\end{equation}
%
Thus Eq.~\eqref{eq:fim} separates sectors by a known collective diagonal
contribution.

This exact-sector effect is the reason the hardware test is scientifically
sharp. If the widened spectral gaps and magnetisation-sector shifts translated
directly into a NISQ protection mechanism, the $\lambda>0$ circuits would be
expected to suppress leakage or improve state retention relative to the
$\lambda=0$ references. The experiments below test precisely that expected
hardware consequence and falsify it for the direct Trotter implementation.

The same identity also explains the digital hardware burden. Because
%
\begin{equation}
  M^2 = nI + 2\sum_{i<j} Z_iZ_j,
  \label{eq:m2_pairwise}
\end{equation}
%
the $\lambda>0$ term is not a single local correction after compilation: it
adds an all-to-all pairwise $Z_iZ_j$ layer, up to a global energy shift. On a
sparse heavy-hex device this all-pair interaction pattern must be routed and
Trotterised, so any ideal spectral gap effect competes directly with extra
two-qubit-gate and depth overhead.

The key scientific boundary follows immediately. Both $H_{XY}$ and
$H_{\mathrm{FIM}}$ commute with total magnetisation, so the ideal Hamiltonian
does not create unitary leakage between magnetisation sectors:
%
\begin{equation}
  [H(\lambda),M]=0.
\end{equation}
%
Therefore any IBM measurement of magnetisation leakage or survival asymmetry is
not a direct ideal-Hamiltonian leakage effect. It must be interpreted as
hardware, noise, state-preparation, transpilation, layout, or readout behaviour.

\section{Artefact-first protocol}

All numerical statements in this manuscript are generated from committed artefacts
under \path{data/scpn_fim_hamiltonian/}. The scripts are:
%
\begin{itemize}
  \item \path{scripts/analyse_fim_spectrum.py},
  \item \path{scripts/analyse_fim_level_spacing.py},
  \item \path{scripts/analyse_fim_entanglement.py},
  \item \path{scripts/analyse_fim_sector_survival.py},
  \item \path{scripts/benchmark_fim_vqe_ground_state.py}, and
  \item \path{scripts/analyse_fim_ibm_repeated_followup.py},
  \item \path{scripts/analyse_fim_readout_matrix_mitigation.py}.
\end{itemize}
%
The claim boundary is maintained separately in
\path{docs/campaigns/scpn_fim_claim_boundary_2026-05-05.md}. This separation is
intentional: paper text is allowed to use only claims that remain valid after
checking the generated artefacts.

\section{Exact spectral effects}

Table~\ref{tab:spectrum} summarises the main exact-spectrum trend. Increasing
$\lambda$ lowers the ground energy and increases the gap and spectral width
for the tested $n=4,6,8$ grid. This is expected from
Eq.~\eqref{eq:sector_shift}: sectors with large $|M|$ are shifted downward.

\begin{table}[tb]
\centering
\small
\begin{tabular}{rrrrr}
\toprule
$n$ & $\lambda$ & $E_0$ & width & gap \\
\midrule
4 & 0 & $-6.3030$ & $12.6060$ & $1.1317$ \\
4 & 4 & $-22.3030$ & $24.6021$ & $12.6060$ \\
6 & 0 & $-10.7930$ & $21.5860$ & $0.4468$ \\
6 & 4 & $-34.7930$ & $40.2628$ & $13.7801$ \\
8 & 0 & $-12.7557$ & $25.2287$ & $0.2827$ \\
8 & 4 & $-44.4730$ & $51.4015$ & $13.7173$ \\
\bottomrule
\end{tabular}
\caption{Exact dense-spectrum summaries generated from
\texttt{fim\_spectrum\_summary\_2026-05-05.json}. The data support a
small-system energy-landscape claim, not a hardware robustness claim.}
\label{tab:spectrum}
\end{table}

\section{Level-spacing and entanglement diagnostics}

The adjacent-gap ratio is used only as a small-system spectral diagnostic. It is
not by itself a proof of strict many-body localisation. For $n=6$, the
full-spectrum mean adjacent-gap ratio changes from $0.4334$ at $\lambda=0$ to
$0.4071$ at $\lambda=4$. For $n=8$, it changes from $0.4368$ at $\lambda=0$ to
$0.4048$ at $\lambda=4$ and $0.3852$ at $\lambda=8$. This supports cautious
language such as ``localisation-like spectral shift'' rather than a definitive
MBL claim.

The low-energy bipartition-entropy artefact shows a related but also bounded
trend. For $n=4$, the mean entropy over the lowest generated eigenstates is
approximately $0.676$ bits for $\lambda\leq 1$, then decreases to
$0.584$ bits at $\lambda=2$, $0.502$ bits at $\lambda=4$, and
$0.490$ bits at $\lambda=8$. The ground state entropy is zero in this grid,
consistent with the large-$|M|$ sector structure selected by the FIM shift.

\section{Ideal sector conservation}

The sector-survival artefact checks the Frobenius norm of the commutator with
$M$ and the off-sector Hamiltonian blocks. For all generated $n=4,6,8$ cases
and all tested $\lambda$ values, the commutator norm with $M$ and the maximum
off-sector block norm are zero in the dense exact representation. The ideal
unitary sector leakage column is therefore exactly zero by construction.

This result is scientifically important because it blocks an incorrect
interpretation. The FIM term can create sector energy barriers and change which
sectors dominate low-energy eigenstates, but it does not create or suppress
ideal magnetisation leakage. A hardware experiment can still be meaningful, but
only as a test of whether the modified circuit family exhibits
$\lambda$-correlated survival or leakage under real noise and controls.

\section{VQE ground-state scoring}

The $n=4$ VQE artefact compares a $K_{ij}$-informed ansatz against
\texttt{TwoLocal} and \texttt{EfficientSU2} at equal repetition count,
three random seeds, and fixed COBYLA iteration budget. The topology-informed
ansatz has lower median relative error at all tested FIM strengths:
%
\begin{table}[tb]
\centering
\small
\begin{tabular}{lrrr}
\toprule
Ansatz & $\lambda=0$ & $\lambda=1$ & $\lambda=4$ \\
\midrule
$K_{ij}$-informed & $0.372\%$ & $0.413\%$ & $0.848\%$ \\
\texttt{TwoLocal} & $7.164\%$ & $11.236\%$ & $9.155\%$ \\
\texttt{EfficientSU2} & $7.708\%$ & $7.228\%$ & $9.050\%$ \\
\bottomrule
\end{tabular}
\caption{Median relative ground-state energy error from
\texttt{fim\_vqe\_ground\_state\_summary\_2026-05-05.json}. This is an
$n=4$ ansatz-scoring result, not a general optimisation theorem.}
\label{tab:vqe}
\end{table}

This result is useful for protocol design because it suggests that the same
topology-informed construction used in the methods paper remains appropriate
for the FIM-augmented Hamiltonian. It does not imply that the ansatz is optimal
at larger $n$ or under hardware noise.

\section{IBM hardware validation}

The IBM hardware component was designed as a falsification test, not as a
demonstration of advantage or error mitigation. Because
$[H(\lambda),M]=0$ in the ideal model, magnetisation leakage on hardware is not
an ideal FIM transition. It is a real-device observable that combines circuit
depth, transpilation, readout, state preparation, relaxation, and coherent gate
errors. The question was therefore narrow: does increasing $\lambda$ produce a
reproducible reduction of measured leakage or increase of state retention under
the tested digital Trotter construction?

The first pilot executed 61 circuits and 249856 shots on \texttt{ibm\_kingston}
as job \texttt{ibm-run-4c0bd60c3fc2c532}. It produced complete raw counts, but it
had one sample per $\lambda$/depth/state condition and therefore supported only
descriptive conclusions. Its direction was already unfavourable to a simple
protection claim: at $\lambda=4$ relative to $\lambda=0$, mean magnetisation
leakage increased by $0.1156$ across 15 matched state/depth comparisons.

We then executed a repeated randomized follow-up on the same backend as job
\texttt{ibm-run-cf4835290f607387}. The repeated design used:
%
\begin{itemize}
  \item $n=4$ qubits,
  \item $\lambda\in\{0,4\}$,
  \item depths $\{2,4,6\}$,
  \item five representative magnetisation-sector initial states,
  \item five randomized replicates per matched condition,
  \item a 16-basis-state readout baseline,
  \item 2048 shots per circuit,
  \item 166 total circuits, and
  \item 339968 total shots.
\end{itemize}

Table~\ref{tab:ibm} summarises the repeated hardware result. The sign is clear
and reproducible in the opposite direction from the protection hypothesis:
$\lambda=4$ increased leakage and reduced exact-state retention relative to
$\lambda=0$ for the tested implementation. The result is therefore a
hardware falsification of this specific implementation, not a positive
coherence result.

\begin{table}[tb]
\centering
\small
\begin{tabular}{lrrrr}
\toprule
Observable & mean $\Delta_{4-0}$ & Fisher $p$ & positive & negative \\
\midrule
State retention & $-0.0797$ & $1.11\times10^{-48}$ & $0/15$ & $15/15$ \\
Magnetisation leakage & $+0.0896$ & $7.49\times10^{-55}$ & $14/15$ & $1/15$ \\
Parity leakage & $+0.0615$ & $6.27\times10^{-48}$ & $13/15$ & $2/15$ \\
Readout-corrected parity leakage & $+0.0644$ & $6.27\times10^{-48}$ & $13/15$ & $2/15$ \\
\bottomrule
\end{tabular}
\caption{Repeated IBM Heron r2 follow-up on \texttt{ibm\_kingston}. Positive
leakage deltas mean $\lambda=4$ leaked more than $\lambda=0$ within matched
initial state and depth. The readout correction is a state-specific parity-flip
correction from the available basis-state readout baselines.}
\label{tab:ibm}
\end{table}

Because the repeated design included all 16 computational-basis readout
baseline circuits, we also applied a literal $2^n\times 2^n$ readout
confusion-matrix inversion offline. The calibration matrix is well-conditioned
for this dataset (condition number $1.049$). The mitigated results preserve the
same sign and scale: state retention changes by $-0.0815$ on average
($p=1.19\times10^{-48}$), magnetisation leakage changes by $+0.0917$
($p=8.14\times10^{-55}$), and parity leakage changes by $+0.0642$
($p=6.19\times10^{-48}$). This removes the scalable-readout-mitigation
objection for the repeated FIM follow-up while leaving the negative hardware
interpretation unchanged.

\section{Discussion}

The generated artefacts support a coherent but limited story. The
Fisher-information-inspired feedback term is not a mysterious source of ideal
sector leakage. It is a collective diagonal term that reshapes the
magnetisation-sector energy landscape. In exact small systems, this reshaping
widens gaps, shifts adjacent-gap diagnostics, and changes low-energy
bipartition entropy. Those are scientifically valid offline findings.

The IBM result answers the tested hardware question negatively. The FIM term
does reshape the exact Hamiltonian spectrum, but in the implemented digital
Trotter circuits the $\lambda=4$ instances carry enough additional hardware
burden that measured leakage increases rather than decreases. This separates
mathematical sector engineering from actual NISQ robustness and prevents an
incorrect coherence-protection claim.

The live-transpilation metadata makes this burden explicit. In the repeated
follow-up artefact, the $\lambda=0$ main circuits averaged transpiled depth
$92.13$ and $24.0$ two-qubit gates, with maxima $134$ and $36$. The
$\lambda=4$ main circuits averaged transpiled depth $354.65$ and
$103.05$ two-qubit gates, with maxima $540$ and $158$. Thus the negative
hardware result is not mysterious: the tested digital implementation asks the
processor to realise an all-to-all $Z_iZ_j$ feedback graph, and the added
noise burden overwhelms the ideal small-system gap widening.

This distinction is the main scientific value of the hardware section. It shows
that a diagonal collective term can be mathematically meaningful while still
being operationally counterproductive after compilation to a noisy
superconducting device. Reporting the negative result is therefore not a failed
experiment; it is the evidence that fixes the claim boundary.

The negative result should not be overgeneralised. It does not prove that all
FIM-like Hamiltonian designs are unhelpful, nor does it exclude other encodings,
native analog implementations, different layouts, lower-depth constructions,
adaptive feedback, or future devices. It does show that the direct all-pair
digital implementation tested here is not a hardware-protection mechanism on
this backend and calibration window.

The next hardware discriminator was a replication-and-stress protocol on two
additional IBM Heron backends. The new lane used the same $n=4$
FIM/baseline observable family with full 16-state readout calibration and
local-folding ZNE scales $\{1,3,5\}$ for the dominant leakage and retention
observables. The result strengthens the falsification rather than reversing
it: on both \texttt{ibm\_marrakesh} and \texttt{ibm\_fez}, $\lambda=4$
continues to increase magnetisation leakage and reduce state retention
relative to $\lambda=0$ after readout mitigation and linear ZNE.

\section{Relation to other protection strategies}

It is useful to distinguish the FIM term from established protection
strategies. Dynamical decoupling applies control pulses designed to average
specific noise components away while preserving the intended logical
evolution. Quantum error correction encodes logical information redundantly
and uses syndrome information to detect or correct errors. The FIM term in
Eq.~\eqref{eq:fim} is neither of these. It is a Hamiltonian design term that
reshapes magnetisation-sector energies, while leaving the ideal
magnetisation-sector decomposition intact.

This distinction explains why the negative IBM result is not surprising in
hindsight. The direct digital implementation adds all-to-all $ZZ$ structure
without adding syndrome extraction, active correction, or a decoupling
sequence targeted to the backend noise spectrum. The correct comparator is
therefore not ``FIM versus no protection'' in the abstract, but ``this
compiled FIM circuit family versus other concrete hardware-control mechanisms
at the same observable and layout boundary.'' Future experiments should compare
the same leakage and retention observables against backend-native
dynamical-decoupling schedules, lower-depth native interaction decompositions,
and explicit QEC or repetition-code scaffolds only after each comparator has
its own preregistered circuit family and resource budget.

To make that comparison falsifiable, the replication/ZNE lane used
$\lambda\in\{0,4\}$, depths $\{2,4\}$, four magnetisation-sector states,
three replicates, local-folding noise scales $\{1,3,5\}$, and full 16-state
readout calibration. The \texttt{ibm\_marrakesh} job
\texttt{d86ubgdg7okc73ell0qg} and the \texttt{ibm\_fez} job
\texttt{d86ubg9789is73906ofg} each executed 160 circuits at 1024 shots. The
readout calibration matrices were well-conditioned, with condition numbers
$1.042$ and $1.077$, respectively. Table~\ref{tab:replicationzne}
summarises the main result.

\begin{table}[tb]
\centering
\small
\begin{tabular}{lrrr}
\toprule
Backend & raw scale-1 & raw linear ZNE & mitigated linear ZNE \\
\midrule
\texttt{ibm\_marrakesh} & $+0.0527$ & $+0.0585$ & $+0.0599$ \\
\texttt{ibm\_fez} & $+0.1897$ & $+0.1862$ & $+0.1937$ \\
\bottomrule
\end{tabular}
\caption{Mean $\lambda=4$ minus $\lambda=0$ magnetisation-leakage deltas
from the dedicated replication/ZNE lane. Positive values mean the FIM
circuits leaked more than the matched baseline circuits. The mitigated
column applies full 16-state readout-matrix inversion before the linear
ZNE fit.}
\label{tab:replicationzne}
\end{table}

The same lane shows retention moving in the opposite direction. After
readout mitigation and linear ZNE, the mean state-retention delta is
$-0.0340$ on \texttt{ibm\_marrakesh} and $-0.1690$ on
\texttt{ibm\_fez}. Thus the replication/ZNE extension does not rescue the
direct digital FIM construction. It makes the claim boundary stronger:
the observed failure is no longer a single-backend Kingston-only artefact,
although it remains IBM-Heron and circuit-family specific.

Finally, we executed larger-width replication/ZNE lanes at $n=6$ and $n=8$ on
the same two Heron backends. These lanes use the same $\lambda=4$ versus
$\lambda=0$ leakage discriminator, local-folding noise scales $\{1,3,5\}$, and
full $2^n$ computational-basis readout calibration. Table~\ref{tab:largerfim}
summarises the mean magnetisation-leakage deltas. The larger-width evidence
does not reverse the negative hardware conclusion: every scale-1 mean delta is
positive, and linear ZNE leaves the deltas positive at both widths and both
backends. The values are backend- and width-dependent, so this is not a
universal leakage law; it is a stronger falsification of the narrow claim that
the tested FIM digital circuit family protects the measured magnetisation
sector on IBM Heron hardware.

\begin{table}[tb]
\centering
\small
\begin{tabular}{llrrr}
\toprule
Backend & $n$ & Raw scale-1 & Raw linear ZNE & Mitigated linear ZNE \\
\midrule
\texttt{ibm\_fez} & 6 & 0.4091 & 0.4021 & 0.4257 \\
\texttt{ibm\_fez} & 8 & 0.1479 & 0.1480 & 0.1539 \\
\texttt{ibm\_marrakesh} & 6 & 0.2474 & 0.2541 & 0.2672 \\
\texttt{ibm\_marrakesh} & 8 & 0.4231 & 0.4203 & 0.4483 \\
\bottomrule
\end{tabular}
\caption{Larger-width FIM replication/ZNE lanes submitted on 2026-05-20.
Values are mean absolute $\lambda=4$ minus $\lambda=0$ magnetisation-leakage
deltas from the retrieved raw-count artefacts. The table addresses the
small-system limitation, but it remains an IBM-Heron circuit-family result
rather than evidence for backend-general protection.}
\label{tab:largerfim}
\end{table}

\section{Limitations}

The study is limited to exact dense diagonalisation at $n=4,6,8$, an
$n=4$ VQE scoring grid, and IBM Heron hardware lanes at $n=4,6,8$. The
level-spacing and entropy results are
localisation-like diagnostics only. They do not prove strict many-body
localisation. The VQE result uses a small optimiser budget and should be treated
as an ansatz-design signal, not a theorem. The IBM hardware result is
multi-backend within IBM Heron devices but still circuit-family specific and
IBM-only. Full readout-matrix mitigation is available for the repeated and
replication FIM follow-ups because they include complete computational-basis
readout baseline
circuits, but this does not correct gate errors, Trotter overhead, layout
routing, or calibration drift. ZNE is also not a universal correction:
it is a stress test for monotone noise-scaling behaviour and can fail when
folded circuits change coherent-error structure, calibration sensitivity, or
dominant leakage channels. Datasets without complete basis calibration remain
limited to exact-state or parity-readout corrections. The Fisher combined
p-values in the hardware table combine matched-condition Welch tests and
should be read as compact evidence of sign stability across the repeated
design, not as proof of backend-general behaviour. No hardware coherence
improvement, quantum advantage, or universal protection claim is made.

At the time of execution, authenticated live QPU access for the hardware
component of this paper was limited to IBM Quantum hardware. The software stack
has provider-neutral execution and artefact-packaging routes for other
gate-model, trapped-ion, neutral-atom, photonic, annealing, analogue, and
simulator targets, but those routes are not evidence of live non-IBM execution.
We plan to rerun the same FIM and baseline payloads on additional hardware
families when compute allocations and provider access become available. Until
those reruns exist, the hardware conclusion is restricted to the IBM Heron
backends, layouts, circuit family, and calibration windows reported here.

\section{Data and code availability}

Code, scripts, and artefacts are maintained in the
\href{https://github.com/anulum/scpn-quantum-control}{project GitHub repository}.
The SCPN/FIM artefacts are under
\path{data/scpn_fim_hamiltonian/}. The manuscript claim boundary is
recorded in
\path{docs/campaigns/scpn_fim_claim_boundary_2026-05-05.md}. The IBM raw-count and
analysis artefacts include:
%
\begin{itemize}
  \item \path{fim_ibm_pilot_raw_counts_2026-05-05_ibm-run-4c0bd60c3fc2c532.json},
  \item \path{fim_ibm_pilot_analysis_2026-05-05_ibm-run-4c0bd60c3fc2c532.json},
  \item \path{fim_ibm_repeated_followup_raw_counts_2026-05-05_ibm-run-cf4835290f607387.json}, and
  \item \path{fim_ibm_repeated_followup_analysis_2026-05-05_ibm-run-cf4835290f607387.json},
  \item \path{fim_readout_matrix_mitigation_summary_2026-05-05_ibm-run-cf4835290f607387.json},
  \item \path{fim_readout_matrix_mitigation_rows_2026-05-05_ibm-run-cf4835290f607387.csv}, and
  \item \path{fim_readout_matrix_mitigation_comparisons_2026-05-05_ibm-run-cf4835290f607387.csv},
  \item \path{fim_replication_zne_submission_ibm_marrakesh_20260520T164759Z.json}, and
  \item \path{fim_replication_zne_submission_ibm_fez_20260520T164759Z.json},
  \item \path{fim_replication_zne_raw_counts_ibm_marrakesh_20260520T165004Z.json},
  \item \path{fim_replication_zne_analysis_ibm_marrakesh_20260520T165004Z.json},
  \item \path{fim_replication_zne_raw_counts_ibm_fez_20260520T165004Z.json}, and
  \item \path{fim_replication_zne_analysis_ibm_fez_20260520T165004Z.json}.
\end{itemize}

\section*{AI usage disclosure}

AI-assisted tools were used for drafting and editing support. Numerical claims,
artefact provenance, scientific framing, authorship, and final responsibility
were verified and retained by the author.

\bibliographystyle{apsrev4-2}
\bibliography{scpn_fim_hamiltonian_refs}

\end{document}
